% preamble-exam-zh.tex — 极简讲解页共享 preamble
% 目标：复刻「题目独立、提示轻框、路线箭头、主体双栏、答案轻收束」的讲义风格。

\documentclass{exam-zh}

% ---------- 基础配置 ----------
\examsetup{
  page/size = a4paper,
  page/show-head = false,
  page/show-foot = false,
  sealline/show = false,
  font = none,
  math-font = none,
}

% ---------- 包 ----------
\usepackage{geometry}
\geometry{top=10mm,bottom=14mm,left=15mm,right=13mm}
\AtBeginDocument{\newgeometry{top=10mm,bottom=14mm,left=15mm,right=13mm}}

\usepackage{xcolor}
\usepackage{needspace}
\usepackage{enumitem}
\usepackage{array}
\usepackage{tabularx}
\usepackage{tcolorbox}
\tcbuselibrary{skins,breakable}
\usepackage{paracol}

% ---------- 打印/屏幕颜色切换 ----------
% 用法：\documentclass[monochrome]{exam-zh} 可降级为灰度。
% 注意：不要使用 \if@xxx 放在 makeatother 外部；这里改成普通条件名。
\newif\ifeduprintcolor
\eduprintcolortrue
\makeatletter
\@ifclasswith{exam-zh}{monochrome}{\eduprintcolorfalse}{}
\makeatother

% ---------- 语义颜色 ----------
\ifeduprintcolor
  \definecolor{edu-blue}{RGB}{31,62,126}
  \definecolor{edu-blue-soft}{RGB}{232,238,250}
  \definecolor{edu-blue-bg}{RGB}{248,250,254}
  \definecolor{edu-ink}{RGB}{30,35,45}
  \definecolor{edu-muted}{RGB}{118,124,136}
  \definecolor{edu-line}{RGB}{209,218,235}
  \definecolor{edu-soft-bg}{RGB}{250,251,253}
  \definecolor{edu-red}{RGB}{168,58,58}
  \definecolor{edu-red-bg}{RGB}{255,247,245}
\else
  \definecolor{edu-blue}{gray}{0.15}
  \definecolor{edu-blue-soft}{gray}{0.90}
  \definecolor{edu-blue-bg}{gray}{0.98}
  \definecolor{edu-ink}{gray}{0.10}
  \definecolor{edu-muted}{gray}{0.45}
  \definecolor{edu-line}{gray}{0.82}
  \definecolor{edu-soft-bg}{gray}{0.97}
  \definecolor{edu-red}{gray}{0.28}
  \definecolor{edu-red-bg}{gray}{0.97}
\fi

% ---------- 全局排版 ----------
\setlength{\parindent}{0pt}
\setlength{\parskip}{0pt}
\linespread{1.16}
\renewcommand{\arraystretch}{1.15}

% ctex 通常提供 \kaishu；这里用它做教师旁白感。
\newcommand{\edukai}{\kaishu}

% ---------- 顶部元信息 ----------
\newcommand{\eduTopBar}[2]{%
  \noindent
  \begin{tabularx}{\linewidth}{@{}Xr@{}}
    {\color{edu-blue}\bfseries #1} &
    {\color{edu-blue}\bfseries #2}
  \end{tabularx}
  \vspace{8mm}
}

% ---------- 轻标题体系 ----------
% 只有真正的大结构才用它；不要给提示/路线滥用标题。
\newcommand{\eduSectionTitle}[1]{%
  \needspace{5\baselineskip}%
  \vspace{2mm}
  {\color{edu-blue}\bfseries\large #1}\par
  \vspace{3mm}
}

\newcommand{\eduSubTitle}[1]{%
  \needspace{4\baselineskip}%
  \vspace{1mm}
  {\color{edu-blue}\bfseries #1}\par
  \vspace{1mm}
}

% ---------- 小问题干：复用 exam (1)(2)(3) 格式 ----------
\newenvironment{eduSubQuestionItem}[1]{%
  \needspace{4\baselineskip}%
  \vspace{2mm}
  \par\noindent
  \hangindent=8mm\hangafter=1
  {\color{edu-blue}\bfseries #1}\enspace
  \begingroup
  \color{edu-ink}
}{%
  \endgroup
  \par
  \vspace{3mm}
}

% ---------- 辅助信息条（解题流程/关键想法/思考）楷体 ----------
\newenvironment{eduAuxItem}[1]{%
  \needspace{4\baselineskip}%
  \vspace{2mm}
  \par\noindent
  \hangindent=8mm\hangafter=1
  {\color{edu-blue}\edukai $\triangleright$~#1}\enspace
  \begingroup
  \color{edu-ink}
  \edukai
}{%
  \endgroup
  \par
  \vspace{4mm}
}

\newcommand{\eduColumnTitle}[2]{%
  {\color{edu-blue}\bfseries #1\hspace{1.5mm}#2}\par
  \vspace{2.5mm}
}

% ---------- 图标：不用外部 icon 字体，保证可移植 ----------
\newcommand{\eduIconBulb}{\(\circ\)}
\newcommand{\eduIconBook}{\(\square\)}
\newcommand{\eduIconCheck}{\(\checkmark\)}
\newcommand{\eduIconPin}{\(\diamond\)}
\newcommand{\eduIconNote}{\(\triangleright\)}

% ---------- 题目区 ----------
% 题目区是页面入口：弱标签 + 一条长蓝线 + 大题干。
\newenvironment{eduProblemBlock}[1][题目]{%
  \needspace{8\baselineskip}
  {\small\color{edu-muted}#1}\par
  \vspace{1.5mm}
  {\color{edu-blue}\hrule height 0.55pt}
  \vspace{4.5mm}
  \begingroup
  \color{edu-ink}
  \large
  \setlength{\parskip}{2.2mm}
  \linespread{1.20}\selectfont
}{%
  \par
  \endgroup
  \vspace{6mm}
}

\newenvironment{eduSubQuestions}{%
  \begin{enumerate}[
    label=(\arabic*),
    leftmargin=8mm,
    itemsep=2.0mm,
    topsep=2.0mm,
    parsep=0pt
  ]
}{%
  \end{enumerate}
}

% ---------- 读题提示：轻框，不做同级标题 ----------
\newtcolorbox{eduReadTipBox}{
  enhanced,
  breakable,
  colback=white,
  colframe=edu-blue!45,
  boxrule=0.45pt,
  arc=1.6mm,
  left=10mm,
  right=4mm,
  top=5mm,
  bottom=5mm,
  before skip=1mm,
  after skip=7mm,
  fontupper=\edukai\normalsize,
  colupper=edu-ink,
  overlay={
    \node[anchor=north west, text=edu-blue] at ([xshift=3.2mm,yshift=-3.2mm]frame.north west)
    {\small\eduIconNote};
  }
}

\newenvironment{eduTipList}{%
  \begin{itemize}[
    label=\textbullet,
    leftmargin=5mm,
    itemsep=1.2mm,
    topsep=0pt,
    parsep=0pt
  ]
}{%
  \end{itemize}
}

% ---------- 解题路线：虚线框 + 文字箭头 ----------
\newenvironment{eduRouteFlow}{%
  \needspace{4\baselineskip}%
  \vspace{1mm}
  \begin{center}
  \normalsize
}{%
  \end{center}
  \vspace{7mm}
}
\newcommand{\eduRouteStep}[1]{%
  {\color{edu-ink}#1}%
}
\newcommand{\eduRouteArrow}{%
  \;$\boldsymbol{\to}$\;%
}

% ---------- 主体讲解：使用 paracol 实现跨页自适应双栏 ----------
\newenvironment{eduDualExplain}{%
  \needspace{6\baselineskip}% 防止标题孤行
  \columnratio{0.25, 0.75}%
  \setlength{\columnsep}{7mm}%
  \columnseprule=0.45pt%
  \colseprulecolor{edu-blue}%
  \begin{paracol}{2}% 开启双栏
}{%
  \end{paracol}% 结束双栏
  \vspace{5mm}
}

\newenvironment{eduExplainLeft}{%
  \normalsize\color{edu-ink}%
}{%
  \switchcolumn
}

\newcommand{\eduColumnDivider}{}

\newenvironment{eduExplainRight}{%
  \normalsize\color{edu-ink}%
}{%
}

\newenvironment{eduNumberList}{%
  \begin{enumerate}[
    label=\arabic*.,
    leftmargin=6mm,
    itemsep=4.8mm,
    topsep=0pt,
    parsep=0pt
  ]
}{%
  \end{enumerate}
}

\newenvironment{eduBulletList}{%
  \begin{itemize}[
    label=\textbullet,
    leftmargin=5mm,
    itemsep=1.2mm,
    topsep=0pt,
    parsep=0pt
  ]
}{%
  \end{itemize}
}

% ---------- 底部双栏：答案 + 方法提醒 ----------
\newenvironment{eduBottomGrid}{%
  \vspace{1mm}
  \par\noindent
}{%
  \par\vspace{1mm}
}

\newenvironment{eduSummaryLeft}{%
  \begin{minipage}[t]{0.30\linewidth}
}{%
  \end{minipage}
}

\newcommand{\eduSummaryDivider}{%
  \hspace{4mm}{\color{edu-line}\vrule width 0.35pt}\hspace{7mm}%
}

\newenvironment{eduSummaryRight}{%
  \begin{minipage}[t]{0.58\linewidth}
}{%
  \end{minipage}
}

% ---------- 答案/方法提醒：轻量收束 ----------
\newtcolorbox{eduSummaryBlock}[2]{
  enhanced,
  breakable,
  colback=white,
  colframe=white,
  boxrule=0pt,
  sharp corners,
  left=0mm,
  right=0mm,
  top=1mm,
  bottom=1mm,
  before skip=1mm,
  after skip=1mm,
  title={\color{edu-blue}\bfseries #1\hspace{1.5mm}#2},
  coltitle=edu-blue,
  attach title to upper={\par\vspace{2mm}},
  fontupper=\normalsize
}

\newenvironment{eduPlainList}{%
  \begin{enumerate}[
    label=(\arabic*),
    leftmargin=7mm,
    itemsep=1.2mm,
    topsep=0pt,
    parsep=0pt
  ]
}{%
  \end{enumerate}
}

% ---------- 兼容旧 block 的轻量盒子 ----------
\newtcolorbox{eduSoftBox}{
  enhanced,
  breakable,
  colback=edu-soft-bg,
  colframe=edu-line,
  boxrule=0.35pt,
  arc=1mm,
  left=3mm,
  right=3mm,
  top=2mm,
  bottom=2mm,
  before skip=1mm,
  after skip=2mm,
  fontupper=\small
}

\newtcolorbox{eduMistakeBox}{
  enhanced,
  breakable,
  colback=edu-red-bg,
  colframe=edu-red!40,
  boxrule=0.35pt,
  arc=1mm,
  left=3mm,
  right=3mm,
  top=2.5mm,
  bottom=2.5mm,
  before skip=1mm,
  after skip=2mm,
  fontupper=\normalsize
}

\newtcolorbox{eduLegacySolution}{
  enhanced,
  breakable,
  colback=edu-soft-bg,
  colframe=edu-line,
  boxrule=0.35pt,
  arc=1mm,
  left=3mm,
  right=3mm,
  top=2mm,
  bottom=2mm,
  before skip=-2mm,
  after skip=4mm,
  fontupper=\small
}

\newtcolorbox{eduTeacherNote}{
  enhanced,
  breakable,
  colback=edu-soft-bg,
  colframe=edu-line,
  boxrule=0pt,
  borderline west={1.5pt}{0pt}{edu-muted},
  left=2.5mm,
  right=2.5mm,
  top=2mm,
  bottom=2mm,
  before skip=1mm,
  after skip=2mm,
  fontupper=\footnotesize,
  colupper=edu-muted
}

% ---------- 答题区 ----------
\newcommand{\answerarea}[1][40mm]{%
  \vspace{3mm}%
  \begin{tcolorbox}[
    enhanced,
    breakable,
    colback=white,
    colframe=edu-line,
    boxrule=0.55pt,
    arc=0.8mm,
    height=#1,
    valign=top,
  ]
  \end{tcolorbox}%
}


\begin{document}

\eduSectionTitle{例题 1 — 求和坐标轴的交点}

\begin{eduProblemBlock}[题目]
直线 $y = -2x + 6$ 与 $x$ 轴、$y$ 轴分别交于点 $A$、$B$。

求 $A$、$B$ 的坐标。
\end{eduProblemBlock}









\begin{eduSubQuestionItem}{求 $B$（$y$ 轴交点）}
求与 $y$ 轴的交点 $B$。\end{eduSubQuestionItem}

\begin{eduDualExplain}
  \begin{eduExplainLeft}
    \eduColumnTitle{\eduIconBulb}{思考引导}
    \begin{eduNumberList}
      \item $y$ 轴上的点横坐标都是 $0$。      \item 令 $x=0$，代入方程 $y = -2x + 6$。    \end{eduNumberList}
  \end{eduExplainLeft}
  \eduColumnDivider
  \begin{eduExplainRight}
    \eduColumnTitle{\eduIconBook}{规范讲解}
    \begin{eduNumberList}
      \item 令 $x = 0$：      \item $y = -2 \times 0 + 6 = 6$      \item $B(0, 6)$    \end{eduNumberList}

  \end{eduExplainRight}
\end{eduDualExplain}




\begin{eduSubQuestionItem}{求 $A$（$x$ 轴交点）}
求与 $x$ 轴的交点 $A$。\end{eduSubQuestionItem}

\begin{eduDualExplain}
  \begin{eduExplainLeft}
    \eduColumnTitle{\eduIconBulb}{思考引导}
    \begin{eduNumberList}
      \item $x$ 轴上的点纵坐标都是 $0$。      \item 令 $y=0$，代入方程解 $x$。      \item 解方程时，$-2x$ 移到右边变 $+2x$。    \end{eduNumberList}
  \end{eduExplainLeft}
  \eduColumnDivider
  \begin{eduExplainRight}
    \eduColumnTitle{\eduIconBook}{规范讲解}
    \begin{eduNumberList}
      \item 令 $y = 0$：      \item $0 = -2x + 6$      \item $2x = 6$      \item $x = 3$      \item $A(3, 0)$    \end{eduNumberList}

  \end{eduExplainRight}
\end{eduDualExplain}




\begin{eduBottomGrid}
  \begin{eduSummaryLeft}
    \begin{eduSummaryBlock}{\eduIconCheck}{答案}
    \begin{eduPlainList}
      \item $A(3, 0)$，$B(0, 6)$    \end{eduPlainList}
    \end{eduSummaryBlock}
  \end{eduSummaryLeft}
  \eduSummaryDivider
  \begin{eduSummaryRight}
    \begin{eduSummaryBlock}{\eduIconPin}{方法提醒}
    \begin{eduBulletList}
      \item $y$ 轴交点：令 $x=0$。      \item $x$ 轴交点：令 $y=0$。      \item 交点坐标写成 $(x, 0)$ 或 $(0, y)$ 的形式。    \end{eduBulletList}
    \end{eduSummaryBlock}
  \end{eduSummaryRight}
\end{eduBottomGrid}



\eduSectionTitle{例题 2 — 交点求k}

\begin{eduProblemBlock}[题目]
直线 $y = kx + 6$ 经过点 $B(-3, 0)$。

(1) 求 $k$ 的值；
(2) 写出这条直线的解析式。
\end{eduProblemBlock}









\begin{eduSubQuestionItem}{(1) 求 $k$}
求 $k$ 的值。\end{eduSubQuestionItem}

\begin{eduDualExplain}
  \begin{eduExplainLeft}
    \eduColumnTitle{\eduIconBulb}{思考引导}
    \begin{eduNumberList}
      \item 把 $x=-3$、$y=0$ 代入 $y=kx+6$。      \item $k \times (-3)$ 等于 $-3k$，注意负号。      \item 移项后两边同时除以 $-3$，负负得正。    \end{eduNumberList}
  \end{eduExplainLeft}
  \eduColumnDivider
  \begin{eduExplainRight}
    \eduColumnTitle{\eduIconBook}{规范讲解}
    \begin{eduNumberList}
      \item 把 $B(-3, 0)$ 代入 $y = kx + 6$：      \item $0 = k \times (-3) + 6$      \item $-3k + 6 = 0$      \item $-3k = -6$      \item $k = 2$    \end{eduNumberList}

  \end{eduExplainRight}
\end{eduDualExplain}




\begin{eduSubQuestionItem}{(2) 写解析式}
写出这条直线的解析式。\end{eduSubQuestionItem}

\begin{eduDualExplain}
  \begin{eduExplainLeft}
    \eduColumnTitle{\eduIconBulb}{思考引导}
    \begin{eduNumberList}
      \item 把求出的 $k=2$ 代回 $y=kx+6$ 即可。    \end{eduNumberList}
  \end{eduExplainLeft}
  \eduColumnDivider
  \begin{eduExplainRight}
    \eduColumnTitle{\eduIconBook}{规范讲解}
    \begin{eduNumberList}
      \item $k = 2$，所以：      \item $y = 2x + 6$    \end{eduNumberList}

  \end{eduExplainRight}
\end{eduDualExplain}




\begin{eduBottomGrid}
  \begin{eduSummaryLeft}
    \begin{eduSummaryBlock}{\eduIconCheck}{答案}
    \begin{eduPlainList}
      \item (1) $k = 2$      \item (2) $y = 2x + 6$    \end{eduPlainList}
    \end{eduSummaryBlock}
  \end{eduSummaryLeft}
  \eduSummaryDivider
  \begin{eduSummaryRight}
    \begin{eduSummaryBlock}{\eduIconPin}{方法提醒}
    \begin{eduBulletList}
      \item 点在直线上 → 坐标代入方程 → 解方程求参数。      \item 负数代入时注意符号。    \end{eduBulletList}
    \end{eduSummaryBlock}
  \end{eduSummaryRight}
\end{eduBottomGrid}



\eduSectionTitle{例题 3 — 两直线交点}

\begin{eduProblemBlock}[题目]
直线 $l_1$：$y = 2x + 1$ 与直线 $l_2$：$y = -x + 4$。

求 $l_1$ 与 $l_2$ 的交点坐标。
\end{eduProblemBlock}









\begin{eduAuxItem}{解题流程}
令 $2x + 1 = -x + 4$，解出 $x$$\;\to\;$把 $x$ 代回任一方程求 $y$$\;\to\;$写出交点坐标 $(x, y)$\end{eduAuxItem}




\begin{eduSubQuestionItem}{求交点}
求 $l_1$ 与 $l_2$ 的交点坐标。\end{eduSubQuestionItem}

\begin{eduDualExplain}
  \begin{eduExplainLeft}
    \eduColumnTitle{\eduIconBulb}{思考引导}
    \begin{eduNumberList}
      \item 交点同时在两条直线上，所以 $y$ 值相等。      \item 令 $2x+1 = -x+4$，把 $-x$ 移到左边变 $+x$。      \item 解出 $x$ 后，代入 $l_1$ 或 $l_2$ 求 $y$（选简单的）。    \end{eduNumberList}
  \end{eduExplainLeft}
  \eduColumnDivider
  \begin{eduExplainRight}
    \eduColumnTitle{\eduIconBook}{规范讲解}
    \begin{eduNumberList}
      \item $2x + 1 = -x + 4$      \item $2x + x = 4 - 1$      \item $3x = 3$      \item $x = 1$      \item 代入 $l_1$：$y = 2 \times 1 + 1 = 3$      \item 交点为 $(1, 3)$    \end{eduNumberList}

  \end{eduExplainRight}
\end{eduDualExplain}




\begin{eduBottomGrid}
  \begin{eduSummaryLeft}
    \begin{eduSummaryBlock}{\eduIconCheck}{答案}
    \begin{eduPlainList}
      \item 交点为 $(1, 3)$    \end{eduPlainList}
    \end{eduSummaryBlock}
  \end{eduSummaryLeft}
  \eduSummaryDivider
  \begin{eduSummaryRight}
    \begin{eduSummaryBlock}{\eduIconPin}{方法提醒}
    \begin{eduBulletList}
      \item 两直线交点 → 联立方程 → 令 $y_1 = y_2$ → 先求 $x$ 再回代求 $y$。      \item 回代时选系数简单的直线，减少出错。    \end{eduBulletList}
    \end{eduSummaryBlock}
  \end{eduSummaryRight}
\end{eduBottomGrid}




\end{document}
